% SPDX-License-Identifier: Apache-2.0
\documentclass[10pt,a4paper]{article}
\usepackage[margin=2.2cm]{geometry}
\usepackage{booktabs}
\usepackage{longtable}
\usepackage{hyperref}
\hypersetup{colorlinks=true,linkcolor=blue,urlcolor=blue}

\title{The Alinx AX7A200B, read out of its own documents}
\author{filmil/a200t\_board\_data}
\date{Built by Bazel from \texttt{data/pins.json}}

\input{counts.tex}

\begin{document}
\maketitle

\begin{abstract}
This document is what the vendor's schematics and user manual say about
one board, the Alinx AX7A200B: a carrier board with an AC7A200 core
board on it, carrying an \texttt{xc7a200tfbg484-2}. Every table here is
generated at build time from \texttt{data/pins.json}, which a tool in
this repository extracts from those documents. Nothing is typed in by
hand, so the document cannot drift from the data, and the data cannot
drift from the documents without the tool being re-run.
\end{abstract}

\section{What the board is}

The board is two boards. The core board, an AC7A200, carries the FPGA,
its DDR3 memory, its configuration flash and its clock. The carrier
carries everything a design talks to: the serial bridge, the LEDs and
buttons, the Ethernet PHY, the HDMI encoder, the PCIe slot, the SFP
cages and the headers. The two meet at a pair of board-to-board
connectors, and a signal crosses them on a net named after the FPGA bank
and the differential pair it comes from, such as \texttt{B15\_L23\_P}.

That naming is why a pin table for this board is work. The core board's
schematic says which ball carries \texttt{B15\_L23\_P}; the carrier's
schematic says what sits on the other end of it; the user manual names
some of those signals and says nothing about others. This repository
joins the three.

The result covers \ballcount{} balls, \namedcount{} of which carry a
named signal, across \bankcount{} I/O banks.

\section{Where the data comes from}

Three documents, each answering a different question.

\begin{itemize}
  \item \emph{AC7A200 Core Board Schematic.pdf} says which FPGA ball
    carries which net, or which of the core board's own signals, and
    what Xilinx calls that ball, such as
    \texttt{IO\_L23P\_T3\_FOE\_B\_15}.
  \item \emph{AX7A200B Carrier Board Schematic.pdf} says which board
    signal sits on each connector net.
  \item \emph{AX7A200B User Manual REV1.0.pdf} gives a table per
    peripheral, with the signal, the ball and a sentence on what the
    signal does.
\end{itemize}

The documents are the vendor's and are not in this repository. What is
here is the table extracted from them, and the tool that extracts it.

\subsection{The trap in the carrier's tables}

The carrier's connector tables cannot be read as text. Each row draws
the net label above and the signal label about three points below it, so
a text extractor that works line by line puts every signal beside the
net of the row \emph{below} it, and every answer is one row out. Read
that way, LED1 takes LED2's ball, \texttt{ETH\_TXCK} moves from P15 to
R16, and the PCIe reset moves from L16 to K16.

The tool therefore reads the carrier as word boxes, with coordinates,
and pairs each signal with the nearest net \emph{above} it. Three facts
check that rule, each from somewhere else:

\begin{enumerate}
  \item the manual's own table says \texttt{ETH\_TXCK} is P15, and the
    rule puts it there;
  \item the rule gives LED1 to L13, LED2 to M13, LED3 to K14 and LED4 to
    K13, which is what the designs in \texttt{HDL/txhdl} drive, and what
    was confirmed by looking at the board;
  \item the rule puts \texttt{UART1\_TXD} on L15, which is the ball
    those designs drive as their serial output on a board whose output
    the host reads.
\end{enumerate}

The tests in this repository assert all three, so a change to the tool
that breaks the rule fails the build.

\section{What to know before using a pin}

\begin{itemize}
  \item \textbf{The serial port's names are written from the FPGA's
    side.} \texttt{UART1\_TXD}, L15, is the FPGA's output and the
    bridge's input.
  \item \textbf{The PCIe reset is L16}, \texttt{IO\_L23P\_T3\_FOE\_B\_15}
    in bank 15, from the slot's pin A11 through a 4.7\,k$\Omega$ pull-up
    to \texttt{+3.3V\_PCIe}. It is a dual-purpose configuration pin,
    \texttt{FOE\_B}, which is free after configuration on a board that
    boots from QSPI.
  \item \textbf{The configuration flash is an N25Q128}, 128\,Mbit, which
    is 16\,MiB. Vivado's flash programming refuses a target that names
    the 256\,Mbit part, saying \texttt{Part identified : mt25ql128 Part
    selected : mt25ql256}.
  \item \textbf{Bank 15 is the 3.3\,V bank} that carries the LEDs, the
    buttons, the serial port and the PCIe reset, so those are
    \texttt{LVCMOS33} without translation.
  \item \textbf{A transceiver ball has two names.} The core board calls
    it \texttt{MGT\_TX0\_P} and the carrier calls it
    \texttt{SFP1\_TX\_P} or \texttt{PCIE\_TX0\_P}. The tables here use
    the functional name and list the other as an alias.
\end{itemize}

\section{The pins}

\input{tables.tex}

\section{Using the data}

\texttt{data/pins.json} is one object per ball:

\begin{verbatim}
{
  "pin": "L16",
  "bank": 15,
  "function": "IO_L23P_T3_FOE_B_15",
  "net": "B15_L23_P",
  "signal": "PCIE_PERST",
  "purpose": "the PCIe slot, x4 in size and x2 in lanes",
  "peripheral": "pcie",
  "board": "carrier",
  "sources": ["AC7A200 Core Board Schematic.pdf",
              "AX7A200B Carrier Board Schematic.pdf"]
}
\end{verbatim}

So a question about the board is a query. Which ball is the PCIe reset:

\begin{verbatim}
jq -r '.pins[] | select(.signal == "PCIE_PERST") | .pin' data/pins.json
\end{verbatim}

What is on bank 16, or which balls the HDMI encoder uses, or what a ball
is for, are the same shape of question.

\end{document}
